% ==============================================================================
% METODOLOGÍA
% ==============================================================================
% Instrucciones de la plantilla:
% - Presentar de manera clara y estructurada el proceso seguido durante la 
%   práctica de laboratorio.
% - Comenzar con la configuración del entorno de trabajo (herramientas, lenguajes,
%   bibliotecas: OpenGL, GLU, Assimp, GLEW, GLFW, GLM, etc.).
% - Detallar los algoritmos gráficos implementados explicando su fundamento teórico
%   y su adaptación práctica (dibujo de formas, transformaciones geométricas, etc.).
% - Incluir imágenes con títulos y descripciones de los pasos ejecutados.
% - Suficientemente detallado para que cualquier lector pueda replicar el programa.
% - REGLA: "Una metodología por actividad".
% ==============================================================================

\section{Metodología}

\subsection{Configuración del entorno de trabajo}
Para el desarrollo de la práctica se utilizó un entorno de desarrollo estructurado en C++ configurado con soporte para OpenGL moderno. En la Tabla \ref{tab:requerimientos} se detallan las herramientas, bibliotecas y componentes del sistema empleados.

\begin{table}[H]
\centering
\renewcommand{\arraystretch}{1.25}
\begin{tabular}{|p{4.2cm}|p{10.5cm}|}
\hline
\textbf{Elemento} & \textbf{Especificación / Versión empleada} \\ \hline
Sistema operativo & Linux / Windows 11 (64 bits) \\ \hline
Entorno de desarrollo & Visual Studio Code / Microsoft Visual Studio \\ \hline
Compilador y estándar & GCC / MSVC (soporte para \texttt{C++20}) \\ \hline
API gráfica & OpenGL 3.3 (Perfil Core) \\ \hline
Bibliotecas principales & GLFW (manejo de ventanas y eventos), GLAD (cargador de funciones), GLM (matemáticas vectoriales y matriciales), Assimp (carga de modelos 3D) \\ \hline
Hardware de prueba & Tarjeta gráfica compatible con OpenGL 3.3 Core \\ \hline
\end{tabular}
\caption{Requerimientos de hardware, software y bibliotecas utilizadas.}
\label{tab:requerimientos}
\end{table}

\subsection{Actividad 1: [Nombre de la Actividad 1]}
Describir detalladamente el procedimiento implementado para la primera actividad: arquitectura del programa, flujo de ejecución y algoritmos gráficos desarrollados.

\subsubsection{Algoritmo y fundamento matemático}
Explicar el fundamento teórico del algoritmo (por ejemplo, el algoritmo de rasterización de líneas de Bresenham o cálculo de transformaciones afines):
\begin{equation}
    P' = M_{proyecci\acute{o}n} \cdot M_{vista} \cdot M_{modelo} \cdot P
    \label{eq:transformacion}
\end{equation}

\subsubsection{Implementación del código}
A continuación se muestra el segmento clave de código que implementa la solución:

\begin{lstlisting}[language=C++, caption={Configuración de buffers (VAO, VBO) y dibujo de geometrías.}, label={lst:actividad1}]
// Generación y enlace del Vertex Array Object (VAO) y Vertex Buffer Object (VBO)
unsigned int VAO, VBO;
glGenVertexArrays(1, &VAO);
glGenBuffers(1, &VBO);

glBindVertexArray(VAO);
glBindBuffer(GL_ARRAY_BUFFER, VBO);
glBufferData(GL_ARRAY_BUFFER, sizeof(vertices), vertices, GL_STATIC_DRAW);

// Definición de atributos de vértices (posición xyz)
glVertexAttribPointer(0, 3, GL_FLOAT, GL_FALSE, 3 * sizeof(float), (void*)0);
glEnableVertexAttribArray(0);
\end{lstlisting}

\subsection{Actividad 2: Aplicación de Modelo de Iluminación y Materiales}
En esta primera actividad se desarrolló la configuración e integración de un modelo de iluminación de Phong utilizando varios objetos geométricos tridimensionales cargado mediante la librería Assimp. El programa establece una arquitectura basada en bucles de renderizado en tiempo real con GLFW y GLAD, gestionando la matriz de vista mediante una cámara interactiva y proyectando un entorno envolvente para contextualizar los reflejos y la iluminación ambiental de la escena.

\subsubsection{Algoritmo y fundamento matemático}
El pipeline gráfico implementa el modelo de iluminación local de Phong, el cual calcula la intensidad de luz reflejada sobre una superficie a partir de tres componentes: ambiental, difusa y especular. Para transformar las coordenadas locales de los vértices del modelo tridimensional al espacio de la pantalla, se aplica la siguiente transformación geométrica fundamental basada en el producto de matrices:
\begin{equation}
    P' = M_{proyecci\acute{o}n} \cdot M_{vista} \cdot M_{modelo} \cdot P
    \label{eq:transformacion}
\end{equation}

Adicionalmente, para poner a los objetos de una apariencia física realista, se configuraron los vectores de coeficientes de material utilizando los valores normalizados de reflectancia obtenidos de la referencia disponible en \url{http://devernay.free.fr/cours/opengl/materials.html}.

\subsubsection{Implementación del código}
A continuación se muestra el segmento clave de código que implementa la solución para la asignación de las propiedades del material y la carga de los shaders de iluminación:

\begin{lstlisting}[language=C++, caption={Configuración de propiedades materiales y renderizado del objeto.}]
// Configuración de propiedades materiales
// Tabla: http://devernay.free.fr/cours/opengl/materials.html
material.ambient = glm::vec4(0.24725f, 0.1995f, 0.0745f, 1.0f);
material.diffuse = glm::vec4(0.75164f, 0.60648f, 0.22648f, 1.0f);
material.specular = glm::vec4(0.628281f, 0.555802f, 0.366065f, 1.0f);
material.transparency = 1.0f;
\end{lstlisting}
